% Source: source/普通高考/2014/2014北京理.pdf, page 1, question 4.
% pdfimages -list found no embedded image objects; the source flowchart is vector artwork.
% Grid evidence: tmp/redraw_flowchart_2014_beijing_science/grid/page1_grid100.jpg and
% q4_block_grid50.jpg; comparison crop is from the no-grid page render.
% Program topology: 开始 -> 输入 m,n 的值 -> k=m,S=1 -> k<m-n+1?;
% 是 -> 输出 S -> 结束; 否 -> S=S*k -> k=k-1 -> return to the shared stem above the test.
% solid segments: all node borders and directed connectors; dashed segments: none.
% Labels are centered in nodes with local parindent=0pt; branch labels sit beside exits.

\begin{tikzpicture}[
    >=Stealth,
    x=1cm,
    y=0.88cm,
    line width=0.45pt,
    line cap=round,
    line join=round,
    font=\normalsize,
    flowtext/.style={anchor=center, align=center,
        execute at begin node={\setlength{\parindent}{0pt}}},
    flowbox/.style={draw, flowtext, minimum height=0.68cm,
        inner xsep=4pt, inner ysep=2pt},
    terminal/.style={flowbox, rounded corners=0.16cm, minimum width=1.45cm},
    io/.style={flowbox, trapezium, trapezium left angle=72, trapezium right angle=108,
        minimum height=0.78cm},
    process/.style={flowbox, minimum width=3.0cm},
    decision/.style={flowbox, diamond, aspect=3.2, minimum width=5.0cm,
        minimum height=1.10cm},
    branchlabel/.style={flowtext, inner sep=0pt},
]
    \path[use as bounding box] (-3.20,1.85) rectangle (5.25,10.35);

    \node[terminal] (start) at (0,9.92) {开始};
    \node[io, minimum width=4.25cm] (input) at (0,8.62) {输入 $m,n$ 的值};
    \node[process, minimum width=3.35cm] (init) at (0,7.30) {$k=m,S=1$};
    \node[decision] (test) at (0,5.45) {$k<m-n+1$};
    \node[io, minimum width=2.15cm] (output) at (0,3.70) {输出 $S$};
    \node[terminal] (finish) at (0,2.35) {结束};

    \node[process, minimum width=2.45cm] (multiply) at (4.00,5.45) {$S=S\cdot k$};
    \node[process, minimum width=2.45cm] (decrement) at (4.00,6.78) {$k=k-1$};

    % Main sequence and affirmative branch.
    \draw[->] (start.south) -- (input.north);
    \draw[->] (input.south) -- (init.north);
    \draw[->] (init.south) -- (test.north);
    \draw[->] (test.south) -- node[branchlabel, left=2pt] {是} (output.north);
    \draw[->] (output.south) -- (finish.north);

    % Negative branch: rightward to multiplication, upward decrement, then a
    % single left-pointing return arrow into the shared stem above the test.
    \draw[->] (test.east) -- (multiply.west);
    \node[branchlabel, anchor=south] at (2.58,5.73) {否};
    \draw[->] (multiply.north) -- (decrement.south);
    \coordinate (loopJoin) at (0,6.78);
    \draw[->] (decrement.west) -- (loopJoin);
\end{tikzpicture}
